\begin{table}[ht]
\caption{\textbf{\psychology example}. The positive document describes the critical equation needed to explain the confusion in the post.}
\centering
% \resizebox{\textwidth}{!}{
\begin{tabular}{p{13cm}}
\toprule
% \hspace{5.5cm} \psychology \\
% \midrule
\textit{\textbf{Query}} \\
\midrule
Why do fNIRS devices commonly use two different frequencies? \\
One of the most common techniques used for functional neuroimaging nowadays is functional near infra-red spectroscopy (fun fact: IIRC Natalie Portman worked on a research paper involving fNIRS as the modality), which shines near infrared light into the brain from a source to a detector (both called optodes) in a "banana" shape. \\
It's not uncommon to read that most of these devices, be they continuous wave (CW) or one of the two kinds that involve fast modulation, frequency domain (FD) or time domain (TD), require two separate frequencies to be emitted. For instance, NIRx explains it as follows on their website: \\
"For neuro-imaging applications it is by far most common to illuminate with two discrete wavelength, which is the minimum requirement to assess relative variations of both oxygenation states of the hemoglobin molecule independently." \\
Why is that the case? \\
I haven't delved into the intricacies of it, but no reason immediately jumps out at me. For instance, in the case of CW, the relative difference in intensity is all that matters, so why do we need two frequencies? \\
\midrule
\textit{\textbf{Chain-of-thought reasoning}} \\
\midrule
In order to understand why we need two frequencies in fNIRS, we will need to check the measurements and calculation in continuous wave systems. 
This may be related to oxygenated (HbO) and deoxygenated (HbR) hemoglobin in the brain. 
We need to refer to the dual-wavelength approach, which is fundamental to fNIRS's ability to provide meaningful information about local brain activity.\\
\midrule
\textit{\textbf{Example positive document}} \\
\midrule
... \\
Where OD is the optical density or attenuation, $I_0$ is emitted light intensity, $I$ is measured light intensity, $\epsilon$ is the attenuation coefficient, $[X]$ is the chromophore concentration, $l$ is the distance between source and detector and $DPF$ is the differential path length factor, and $G$ is a geometric factor associated with scattering. \\
When the attenuation coefficients $\epsilon$ are known, constant scattering loss is assumed, and the measurements are treated differentially in time, the equation reduces to: \\
$\Delta [X] = \Delta OD/(\epsilon d)$ \\
Where d is the total corrected photon path-length. \\
Using a dual wavelength system, measurements for HbO2 and Hb can be solved from the matrix equation: \\
$
\begin{pmatrix}
\Delta OD_{\lambda 1} \\
\Delta OD_{\lambda 2}
\end{pmatrix}
$
=
$
\begin{pmatrix}
\epsilon^{Hb}_{\lambda 1} d & \epsilon^{HbO_2}_{\lambda 1} d \\
\epsilon^{Hb}_{\lambda 2} d & \epsilon^{HbO_2}_{\lambda 2} d 
\end{pmatrix}
$
$
\begin{pmatrix}
\Delta [X]^{Hb} \\
\Delta [X]^{HbO_2}
\end{pmatrix}
$ \\
... \\
\midrule
\textit{\textbf{Example negative document}} \\
\midrule
What Is Wave Frequency? \\
The number of waves that pass a fixed point in a given amount of time is wave frequency. Wave frequency can be measured by counting the number of crests (high points) of waves that pass the fixed point in 1 second or some other time period. The higher the number is, the greater the frequency of the waves. The SI unit for wave frequency is the hertz (Hz), where 1 hertz equals 1 wave passing a fixed point in 1 second. The Figure below shows high-frequency and low-frequency transverse waves. \\
Q: The wavelength of a wave is the distance between corresponding points on adjacent waves. For example, it is the distance between two adjacent crests in the transverse waves in the diagram. Infer how wave frequency is related to wavelength. \\
A: Waves with a higher frequency have crests that are closer together, so higher frequency waves have shorter wavelengths. \\
% Wave Frequency and Energy \\
% The frequency of a wave is the same as the frequency of the vibrations that caused the wave. For example, to generate a higher-frequency wave in a rope, you must move the rope up and down more quickly. This takes more energy, so a higher-frequency wave has more energy than a lower-frequency wave with the same amplitude. You can see examples of different frequencies in the Figure below (Amplitude is the distance that particles of the medium move when the energy of a wave passes through them.) \\
... \\
\bottomrule
\end{tabular}
% }
\label{tab:psychology_example}
\end{table}